\clearpage
\chapter{节奏、乐句与曲式进阶}
\label{chap:rhythm-phrase-form}

\subsection*{学完这一章能做什么}
\begin{itemize}
  \item 分清 6/8 与 3/4 的听觉差别，能写出 hemiola（三对二）而不把拍号写乱。
  \item 用“动机—发展—乐句—终止”拆一段八小节，指出哪一拍是句眼、哪一拍收束。
  \item 给一段 16 小节结构标出 verse/chorus/bridge，说明每段的节奏功能而不是只写标签。
  \item 在夜半里用同一动机做四次变化（模进、倒影、扩大、碎片化），再拼成八小节。
\end{itemize}

\subsection{6/8 不是“每小节六个四拍”}

3/4 是三个等重拍，6/8 是“三拍一组的复合拍”：六个八分音符分成两大组，每组三个。听感上 3/4 更像“圆舞曲”，6/8 更像“摇摆/船歌”。写作时先决定“大拍”在哪：

\begin{center}
\small
\begin{tabular}{clll}
\toprule
拍号 & 大拍 & 细分 & 常见用途 \\
\midrule
3/4 & 3 个四分拍 & 每拍二分/四分 & 华尔兹、进行曲中的三拍 \\
6/8 & 2 个大拍（每组 3 个八分） & 每大拍三个八分 & 船歌、抒情副歌、摇摆 \\
12/8 & 4 个大拍（每组 3 个八分） & 慢蓝调、福音 & 每大拍三连音感 \\
\bottomrule
\end{tabular}
\end{center}

hemiola 是在同一段里把“两组三拍”临时听成“三组两拍”（或反过来）。6/8 中把三个大组各切两半，会产生 3/4 的错觉；3/4 中把两组各拉成三拍，会产生 6/8 的错觉。hemiola 适合放在乐句收尾，制造“要转段”的悬念。

\begin{figure}[htbp]
\centering
\begin{tikzpicture}[x=0.75cm,y=0.5cm,font=\scriptsize]
  \foreach \b in {0,...,11} { \draw[gray] (\b,0)--(\b,2.2); }
  \foreach \b/\lab in {0/1,1/2,2/3,3/4,4/5,5/6} { \node at (\b+0.5,2.35) {\lab}; }
  \draw[thick] (0,0)--(3,0)--(3,0.8)--(6,0.8);
  \draw[thick] (6,0)--(9,0)--(9,0.8)--(12,0.8);
  \node at (3,1.3) {大拍 1};
  \node at (9,1.3) {大拍 2};
  \node[gray] at (1.5,-0.4) {三个八分};
  \node[gray] at (7.5,-0.4) {三个八分};
\end{tikzpicture}
\caption*{图 12-1：6/8 的两个大拍。每个大拍内部有三个八分；写旋律时不要把六个八分平均当成六个独立拍。}
\end{figure}

\subsection{律动：重音、切分与三对二}

律动（groove）不是“节奏型名称”，而是重音周期与音色、时值的组合。三条可操作规则：

\begin{enumerate}
  \item \textbf{重音搬家}：把本应落在强拍的重音提前或推迟半拍，制造切分；切分后下一拍要留一个“落地”音，否则会飘。
  \item \textbf{三对二}：在 4/4 里让一条声部按三连音、另一条按八分，形成 3:2；两条都要有稳定周期，不能只是“乱”。
  \item \textbf{留白}：连续切分后留半拍空拍，让听众的耳朵重新对齐；满拍切分会让段落失去方向。
\end{enumerate}

写作时先画“重音图”：每小节标出强、弱、次强、弱，再把旋律重音与鼓组重音对比。若所有重音都落在同一拍，律动会平；若每拍都切分，律动会散。

\subsection{乐句：动机、句眼与终止}

乐句（phrase）是一个能独立呼吸的旋律单位，通常 2–4 小节。分析时找三样：

\begin{itemize}
  \item \textbf{动机}（motif）：最短的可识别音型，通常 2–5 个音；
  \item \textbf{句眼}：乐句里张力最高或信息最密的一拍，常是模进的顶点；
  \item \textbf{终止}：乐句怎样收：半终止（停在 V）、正格终止（V--I）、变格终止（IV--I）、阻碍终止（V--vi）。
\end{itemize}

动机发展的四种基本手法：

\begin{center}
\small
\begin{tabular}{cll}
\toprule
手法 & 做法 & 听觉效果 \\
\midrule
模进 sequence & 同一音型整体上/下移 & 推进、可预期 \\
倒影 inversion & 音程方向反过来 & 熟悉又陌生 \\
扩大 augmentation & 时值整体拉长 & 强调、收束 \\
碎片化 fragmentation & 只取动机的一半 & 加速、紧张 \\
\bottomrule
\end{tabular}
\end{center}

\subsection{曲式：段落的功能，不是标签}

常见 16 小节结构可写作 AABA、verse--chorus、intro--verse--chorus--bridge--outro。标签本身不重要，重要的是每段的节奏/和声功能：

\begin{itemize}
  \item \textbf{主歌 verse}：叙事、低张力，节奏较稳，和声多用 I--vi--IV--V 或循环；
  \item \textbf{副歌 chorus}：记忆点，节奏更密或更高，和声常回到 I，句尾用强终止；
  \item \textbf{桥段 bridge}：对比，节奏可突然变慢/变快，和声常用借用或转调；
  \item \textbf{过渡 transition}：用 hemiola、切分或上升模进把耳朵送到下一段。
\end{itemize}

一段八小节可以按“2+2+4”拆：前两小节立动机，第三四小节模进/倒影，后四小节做句眼与终止。若八小节全是同一节奏型，听感会像 loop；若每两小节换一次拍号，听感会像练习曲。

\subsection{练习}
\begin{enumerate}
  \item 6/8 与 3/4 的都是六个八分，为什么听感不同？\\
  \emph{答案：}6/8 分成两个大拍（每组三个八分），3/4 是三个等重四分拍；重音周期不同。
  \item 写一个 4/4 的两小节节奏，第二小节用 hemiola 暗示 3/4。\\
  \emph{答案示例：}第一小节四拍均分；第二小节把 4 个四分写成 3+3+2 的八分分组。
  \item 给“1 2 3 5”这个四音动机做模进、倒影、扩大各一次。\\
  \emph{答案示例：}模进 2 3 4 6；倒影 1 7 6 4（以 1 为轴）；扩大 1 1 2 2 3 3 5 5（时值翻倍）。
  \item 分析一段八小节，标出句眼与终止类型。\\
  \emph{答案：}句眼通常是模进顶点或最高音；终止看最后两小节低音与和弦，V--I 为正格，V--vi 为阻碍。
  \item 为什么连续切分后要留白？\\
  \emph{答案：}让听众重新对齐拍点；没有落地的切分会失去方向感。
  \item 在夜半里用“1 2 3 5”做模进、倒影、扩大、碎片化四次变化，拼成八小节；至少一次变化要换拍号或加 hemiola。\\
  \emph{答案示例：}1--2 小节原形，3--4 小节模进，5--6 小节倒影，7--8 小节扩大＋碎片化；第 8 小节用 3+3+2 的 hemiola 收束。
\end{enumerate}

\subsection{帮助主题映射}
\begin{center}
\footnotesize
\setlength{\tabcolsep}{3pt}
\begin{longtable}{p{0.34\linewidth}p{0.45\linewidth}p{0.15\linewidth}}
\toprule
本章小节 & 应用内帮助锚点 & 备注 \\
\midrule
\endhead
6/8 与 hemiola & \texttt{music-theory-compound-meter} & 图 12-1 \\
重音与三对二 & \texttt{music-theory-groove-hemiola} & 规则 \\
乐句与动机 & \texttt{music-theory-phrase-motif} & 表 12-1 \\
段落功能 & \texttt{music-theory-section-function} & 结构 \\
在夜半里怎么做 & \texttt{music-theory-motif-variations} & 操作 \\
练习 & \texttt{music-theory-ch12-practice} & 答案在本章 \\
\bottomrule
\end{longtable}
\end{center}
